%% Conference-style paper — Fashion-MNIST augmentation scaling study
\documentclass[11pt]{article}

%% ── Packages ──────────────────────────────────────────────────────────────────
\usepackage[T1]{fontenc}
\usepackage[utf8]{inputenc}
\usepackage{lmodern}
\usepackage[a4paper, margin=1in]{geometry}
\usepackage{natbib}
\usepackage{graphicx}
\usepackage{booktabs}
\usepackage{amsmath}
\usepackage{microtype}
\usepackage{hyperref}
\usepackage{xcolor}
\usepackage{caption}
\usepackage{subcaption}
\usepackage{parskip}
\usepackage{enumitem}

\hypersetup{
  colorlinks=true,
  linkcolor=blue!60!black,
  citecolor=green!50!black,
  urlcolor=blue!70!black,
}

%% ── Title block ───────────────────────────────────────────────────────────────
\title{\textbf{Diminishing Harm, Not Emerging Benefit: Basic Geometric Augmentation
Consistently Hurts a Small CNN on Fashion-MNIST
Within a 5-Epoch Budget}}

\author{%
  OpenScientist Report Agent\\[2pt]
  \small\textit{Automated Empirical Study --- 2026-10-09}
}
\date{}

%% ══════════════════════════════════════════════════════════════════════════════
\begin{document}
\maketitle
\thispagestyle{empty}

%% ── Abstract ──────────────────────────────────────────────────────────────────
\begin{abstract}
We present a controlled empirical study of how basic geometric data augmentation
(random horizontal flip and affine transform) affects test accuracy for a small
convolutional neural network (CNN) trained on Fashion-MNIST across seven training
set sizes ranging from 500 to 60,000 samples, under a strict 5-epoch training
budget.
Across 28 training runs (7 sizes~$\times$ 2 augmentation conditions~$\times$
2 random seeds) executed on an NVIDIA~L4 GPU with PyTorch~2.11.0, augmentation
\emph{consistently decreased} test accuracy at every tested size: mean
$\Delta = \text{acc}_\text{aug} - \text{acc}_\text{noaug}$ ranges from
$-8.06$~pp at $N{=}500$ to $-2.81$~pp at $N{=}60{,}000$.
Both seeds independently agree on the negative sign at all seven sizes
(concordance fraction 1.00), ruling out seed artefacts.
A log-linear OLS regression confirms that augmentation harm decreases with
training set size ($b{=}{+}0.530$, $R^2{=}0.613$, $p{=}0.037$), but the
estimated crossover to positive benefit lies at $N^\ast \approx 6.45$ million
samples---far outside the tested range---with a bootstrap 95\% confidence
interval spanning four orders of magnitude ($[40{,}325;\, {\sim}10^9]$).
The mechanism is \emph{convergence overhead}: augmented runs exhibit higher
epoch-1 training loss in 14/14 configurations, and the 5-epoch budget is
insufficient for the model to recover at any tested size.
The no-augmentation baseline at $N{=}60{,}000$ achieves 91.48\% test accuracy,
confirming correct CNN calibration.
These results constitute cautionary evidence that short training budgets can
invert augmentation's expected benefit, independent of dataset size.
\end{abstract}

%% ── 1. Introduction ───────────────────────────────────────────────────────────
\section{Introduction}
\label{sec:intro}

Data augmentation is a widely used regularisation strategy for training CNNs,
widely assumed to be beneficial for generalisation, and especially valuable in
low-data regimes where the model would otherwise overfit
\citep{mukhamediev2024state,kumar2025enhancing}.
Canonical geometric transforms---horizontal flipping, rotation, translation,
and cropping---are standard practice across benchmark datasets
\citep{karada2020empirical,kang2024game}.

However, augmentation's benefit is not unconditional.
\citet{karada2020empirical} demonstrated on CIFAR-10 that the magnitude of
augmentation benefit is sensitive to the size of the available training set,
suggesting a nuanced interaction between sample count and augmentation
effectiveness.
Building on this precedent, we ask a more refined question: \emph{does basic
geometric augmentation provide a positive benefit to a small CNN on
Fashion-MNIST, and how does this benefit scale with training set size from
500 to 60,000 samples under a fixed 5-epoch training budget?}

Our study tests the hypothesis that augmentation benefit is largest at small
$N$ (500--5,000 samples) and diminishes to negligibility around
$N \approx 10{,}000$--$20{,}000$ samples, at which point the training set is
large enough that augmentation is no longer necessary.
We find the hypothesis is not supported.
Instead, augmentation was \emph{not beneficial} under the 5-epoch budget at any
tested size; harm decreases as $N$ grows, but no crossover to positive benefit
is observed within the tested regime.

\textbf{Contributions.}
(1)~We provide the first systematic quantification of augmentation-benefit
scaling on Fashion-MNIST across a wide range of training set sizes under a
fixed architecture and epoch count.
(2)~We document a cautionary negative result: basic geometric augmentation
consistently hurts a small CNN under a 5-epoch training budget across all
tested sizes.
(3)~We identify \emph{convergence overhead}---higher per-step training loss
induced by augmentation---as the mechanistic explanation, and show it persists
from $N{=}500$ to $N{=}60{,}000$.

%% ── 2. Related Work ───────────────────────────────────────────────────────────
\section{Related Work}
\label{sec:related}

\paragraph{Small CNN benchmarks on Fashion-MNIST.}
Small CNN architectures routinely achieve 88--94\% accuracy on the full
Fashion-MNIST training set \citep{mukhamediev2024state,yee2026based}.
Our no-augmentation baseline of 91.48\% at $N{=}60{,}000$ falls squarely in
this range, validating the experimental setup.
Several studies have examined architecture choices (depth, batch normalisation,
dropout) and hyperparameter effects on Fashion-MNIST
\citep{meshkini2019analysis,xhaferra2022classification}, but none systematically
vary training set size across a controlled sample-count range.

\paragraph{Augmentation benefit and dataset size.}
The relationship between augmentation benefit and dataset size has been explored
most thoroughly on CIFAR-10.
\citet{karada2020empirical} showed that augmentation benefit varies with training
set size, motivating the replication and extension of this analysis to
Fashion-MNIST.
\citet{kang2024game} demonstrated augmentation benefits with MFG Augment, a
Mean-Field Game path-interpolation method that synthesises continuous image
distribution transitions, which is substantially more sophisticated than the
simple geometric transforms used here.

\paragraph{Augmentation in low-data domains.}
In speciality domains, augmentation has been widely reported to improve
performance in low-data regimes.
\citet{pham2018deep} found augmentation helpful for skin-lesion classification,
and \citet{harun2024dysgraphia} reported similar findings for dysgraphia
detection; both studies used training runs of 50 or more epochs.
\citet{snider2023using} showed augmentation benefits ultrasound image
classification, again with substantially longer training than 5 epochs.
The apparent discrepancy between those findings and ours is, as we demonstrate,
fully explained by the epoch budget: 5 epochs are insufficient for a CNN to
overcome augmentation's convergence overhead regardless of training set size.

\paragraph{Augmentation transform design.}
\citet{poojary2021effect} and \citet{zhang2020circular} examined the effect of
specific augmentation types on image classification accuracy, highlighting
that not all transforms are equally beneficial.
\citet{kumar2025enhancing} proposed three novel augmentation techniques---pairwise
channel transfer, object occlusion, and masking---that improve CNN performance
on small datasets, demonstrating that carefully designed augmentation strategies
can meaningfully reduce overfitting in limited-data regimes.

%% ── 3. Methods ────────────────────────────────────────────────────────────────
\section{Methods}
\label{sec:methods}

\subsection{Dataset and Subsets}
Fashion-MNIST \citep{mukhamediev2024state} comprises 60,000 training and 10,000
test grey-scale $28{\times}28$ images across 10 balanced clothing categories.
We created seven stratified class-balanced training subsets of sizes
$N \in \{500,\, 1{,}000,\, 2{,}500,\, 5{,}000,\, 10{,}000,\, 30{,}000,\, 60{,}000\}$
by sampling uniformly within each class using random seeds 42 and 123.
Each subset was further split 90/10 into a training partition and a validation
partition used for monitoring only (no early stopping).
All conditions were evaluated on the fixed 10,000-sample test set.

\subsection{Architecture}
The CNN consists of two dual-convolution blocks followed by a fully connected
head:
\begin{itemize}[nosep]
  \item \textbf{Block 1}: Conv($1{\to}32$) $\to$ BN $\to$ ReLU $\to$
        Conv($32{\to}32$) $\to$ BN $\to$ ReLU $\to$ MaxPool($2{\times}2$)
        $\to$ Dropout(0.25)
  \item \textbf{Block 2}: Conv($32{\to}64$) $\to$ BN $\to$ ReLU $\to$
        Conv($64{\to}64$) $\to$ BN $\to$ ReLU $\to$ MaxPool($2{\times}2$)
        $\to$ Dropout(0.25)
  \item \textbf{FC head}: Linear($3{,}136{\to}256$) $\to$ ReLU
        $\to$ Dropout(0.50) $\to$ Linear($256{\to}10$)
\end{itemize}
The total parameter count is 871,018, confirmed before training.
All kernel sizes are $3{\times}3$ with padding~1.

\subsection{Augmentation Policy}
The \textit{aug} condition applies, during training only,
\texttt{RandomHorizontalFlip}$(p{=}0.5)$ followed by
\texttt{RandomAffine}(degrees$=\pm10^\circ$, translate$=(0.1,\,0.1)$).
The \textit{noaug} condition applies no training-time transforms.
In both conditions, test-time preprocessing is identical: normalise with
Fashion-MNIST channel mean 0.2860 and standard deviation 0.3530.

\subsection{Training Protocol}
All runs use Adam (learning rate $10^{-3}$), batch size 64,
cross-entropy loss, and exactly 5 epochs.
No learning-rate scheduler is used.
Weights are initialised using PyTorch defaults.
Each (size, condition, seed) triple constitutes one independent run;
random seeds 42 and 123 are applied to Python's \texttt{random}, NumPy, and
PyTorch (including CUDA) at the start of each run.
All 28 runs were executed sequentially on an NVIDIA~L4 GPU under
PyTorch~2.11.0+cu130.

\subsection{Analysis}
For each of the 7 training set sizes we compute the per-seed augmentation
benefit $\Delta_s = \text{acc}_{\text{aug},s} - \text{acc}_{\text{noaug},s}$
and the mean $\bar\Delta$ across seeds $s \in \{42, 123\}$.
We fit an OLS log-linear regression
$\bar\Delta = a + b\cdot\log_2(N)$
and report the slope $b$, coefficient of determination $R^2$, and a
$t$-distribution 95\% confidence interval for $b$.
The crossover $N^\ast$ (where $\bar\Delta = 0$) is estimated as
$N^\ast = 2^{-a/b}$ and its uncertainty is quantified via 5,000 bootstrap
iterations (percentile 95\% CI, bootstrap seed~42, right-tail values clipped
at $10^9$).
Spearman's $\rho$ is computed for monotonicity assessment.
All analysis is purely descriptive given $n{=}2$ seeds per cell.

%% ── 4. Results ────────────────────────────────────────────────────────────────
\section{Results}
\label{sec:results}

\subsection{Augmentation Consistently Hurts at All Training Set Sizes}

Table~\ref{tab:results} and Figure~\ref{fig:delta} present the primary results.
The mean augmentation benefit $\bar\Delta$ is negative at every tested training
set size, ranging from $-8.06$~pp at $N{=}500$ to $-2.81$~pp at $N{=}60{,}000$.
Both random seeds independently produce negative $\Delta$ at all 7 sizes
(concordance fraction 1.00), confirming the result is not a seed artefact.
Figure~\ref{fig:bars} illustrates that no-augmentation accuracy exceeds
augmented accuracy at every size; the gap narrows with $N$ but never closes.

\input{table_results.tex}

\begin{figure}[t]
  \centering
  \includegraphics[width=0.82\linewidth]{../figures/fig_delta_vs_size.png}
  \caption{Augmentation benefit $\Delta = \text{acc}_\text{aug} -
  \text{acc}_\text{noaug}$ (percentage points) versus training set size $N$
  (log$_2$ scale) for a small CNN on Fashion-MNIST trained for 5 epochs.
  Grey markers show individual seed values (circles: seed~42; squares: seed~123);
  the black connected line shows mean $\bar\Delta$;
  the blue line is the OLS log-linear fit
  $\hat\Delta = -12.00 + 0.530\cdot\log_2(N)$
  ($R^2{=}0.613$, $p{=}0.037$) with ${\pm}1.96\,\text{SE}$ confidence band.
  All 7 $\bar\Delta$ values are negative.
  The OLS-estimated crossover $N^\ast \approx 6.45$ million
  (bootstrap 95\% CI $[40\text{K};\, {\sim}10^9]$) lies far beyond the tested
  range and is shown for illustration only.}
  \label{fig:delta}
\end{figure}

\begin{figure}[t]
  \centering
  \includegraphics[width=0.82\linewidth]{../figures/fig_test_acc_bars.png}
  \caption{Test accuracy at epoch~5 (\%) for no-augmentation (blue) and
  augmented (orange) conditions, grouped by training set size $N$.
  Individual seed values are overlaid as scatter dots (circles: seed~42;
  squares: seed~123).
  No-augmentation accuracy exceeds augmented accuracy at every size;
  the gap narrows from approximately 8~pp at $N{=}500$ to approximately 2.8~pp
  at $N{=}60{,}000$ but remains substantial throughout.}
  \label{fig:bars}
\end{figure}

\subsection{Augmentation Harm Decreases with Training Set Size}

The log-linear OLS fit yields slope $b{=}{+}0.530$ (95\% CI $[0.046,\, 1.015]$),
$R^2{=}0.613$, $p{=}0.037$, confirming that augmentation harm decreases
significantly as training set size grows on a log scale.
The Spearman rank correlation is $\rho{=}0.714$ ($p{=}0.071$, $n{=}7$), a
positive rank correlation consistent with the trend but not reaching conventional
significance given the small number of size levels.

Notably, the trend is non-monotone: mean $\bar\Delta$ is $-4.77$~pp at
$N{=}2{,}500$ but worsens to $-6.55$~pp at $N{=}5{,}000$ and remains at
$-6.14$~pp at $N{=}10{,}000$ before recovering at larger sizes.
This mid-range dip is not captured by the log-linear model (which contributes to
$R^2{=}0.613$ rather than a higher value) and is discussed further in
Section~\ref{sec:discussion}.

\subsection{The Estimated Crossover is Far Outside the Tested Range}

The OLS-implied crossover $N^\ast$ (where $\bar\Delta{=}0$) is approximately
$6.45$ million samples.
This estimate is flagged as extrapolated: it lies more than two orders of
magnitude beyond the largest tested training set ($N{=}60{,}000$).
The bootstrap 95\% CI spans four orders of magnitude
($[40{,}325;\, {\sim}10^9]$), providing negligible practical predictive value.
This result should be treated as an illustration of the regression framework, not
a reliable estimate of an actionable threshold.

\subsection{Convergence Overhead as the Mechanistic Explanation}

Figure~\ref{fig:curves_all} shows training loss and validation accuracy curves for
all 7 sizes.
Augmented runs exhibit higher epoch-1 training loss in \emph{all} 14
(seed~$\times$~size) configurations.
Representative examples: at $N{=}500$ (seed~42), epoch-1 loss is 2.467
(augmented) vs.\ 2.327 (no-augmentation); at $N{=}60{,}000$ (seed~42),
the corresponding values are 0.772 vs.\ 0.506
(Figures~\ref{fig:curves500}--\ref{fig:curves60k}).
In every case, augmented runs begin with a harder optimisation landscape and do
not recover to match the no-augmentation baseline within 5 epochs.
This \emph{convergence overhead}---harder per-step examples that reduce per-step
gradient informativeness---is the principal mechanism behind the universally
negative $\bar\Delta$.

\begin{figure}[p]
  \centering
  \includegraphics[width=\linewidth]{../figures/fig_curves_all.png}
  \caption{Per-epoch training loss (left panels) and validation accuracy
  (right panels) for all 7 training set sizes (rows, $N{=}500$ to $60{,}000$)
  under no-augmentation (blue) and augmented (orange) conditions.
  Individual seeds are shown as dashed lines; means as solid lines.
  Augmented runs begin with consistently higher epoch-1 loss and converge more
  slowly, explaining why 5 epochs are insufficient for augmentation to provide
  net benefit at any tested size.}
  \label{fig:curves_all}
\end{figure}

\begin{figure}[htbp]
  \centering
  \begin{subfigure}[b]{0.48\linewidth}
    \includegraphics[width=\linewidth]{../figures/curves_size500.png}
    \caption{$N{=}500$.}
    \label{fig:curves500}
  \end{subfigure}
  \hfill
  \begin{subfigure}[b]{0.48\linewidth}
    \includegraphics[width=\linewidth]{../figures/curves_size60000.png}
    \caption{$N{=}60{,}000$.}
    \label{fig:curves60k}
  \end{subfigure}
  \caption{Training loss and validation accuracy curves at the two extreme
  training set sizes.
  \textbf{(a)} At $N{=}500$, augmented runs begin with substantially higher
  epoch-1 loss (${\approx}2.47$ vs.\ ${\approx}2.33$ mean) and do not recover
  within 5 epochs.
  \textbf{(b)} At $N{=}60{,}000$, augmented runs still start with approximately
  50\% higher epoch-1 loss (seed~42: 0.772 vs.\ 0.506) and trail
  no-augmentation by ${\sim}2.8$~pp at epoch~5, illustrating that the
  5-epoch budget is insufficient for augmentation to provide net benefit even
  with the complete Fashion-MNIST training set.}
  \label{fig:curves_extremes}
\end{figure}

\subsection{Baseline Calibration}

The no-augmentation baseline at $N{=}60{,}000$ achieves 91.40\% (seed~42) and
91.57\% (seed~123), with a mean of 91.48\%.
This falls within the 88--92\% range expected for small CNNs on Fashion-MNIST
\citep{mukhamediev2024state}, confirming that the architecture and training
protocol are well-calibrated.

%% ── 5. Discussion ─────────────────────────────────────────────────────────────
\section{Discussion}
\label{sec:discussion}

\paragraph{Convergence overhead explains universal negative $\bar\Delta$.}
The 14/14 higher epoch-1 loss pattern under augmentation is the central
mechanistic finding.
By randomly flipping and affinely distorting each image, augmentation increases
the effective difficulty of each training example.
For a 5-epoch budget, the model never recovers the per-step gradient
informativeness it sacrifices to handle harder examples; it is effectively
under-trained relative to the no-augmentation condition at every tested size.
This interpretation is consistent with the observation that $|\bar\Delta|$
decreases as $N$ grows: with larger training sets, more gradient steps are
available per epoch (given fixed batch size), partially compensating for the
increased per-step difficulty.

\paragraph{Non-monotone dip at $N{=}5{,}000$--$10{,}000$.}
The worsening of $\bar\Delta$ from $-4.77$~pp at $N{=}2{,}500$ to
$-6.55$~pp at $N{=}5{,}000$, followed by gradual recovery at larger sizes, is a
genuine irregularity that the log-linear model only partially captures.
One plausible explanation is that at $N{\approx}2{,}500$--$5{,}000$ the
no-augmentation condition begins to overfit less aggressively (as more training
data becomes available), while the augmented condition has not yet accumulated
enough gradient steps to compensate for convergence overhead; this temporarily
widens the gap before more training data at $N{>}10{,}000$ allows partial
recovery.
We treat this dip as an observation to be reproduced and explained in future
work rather than as an artifact of the current setup.

\paragraph{Comparison with prior work.}
Studies finding augmentation beneficial
\citep{pham2018deep,snider2023using,kumar2025enhancing} typically employ
50--200+ training epochs.
The apparent discrepancy with our finding is fully explained by epoch budget
differences; we make no claim that augmentation is inherently inferior to
no-augmentation for Fashion-MNIST, only that it is not beneficial within a
5-epoch budget.
\citet{karada2020empirical} found dataset-size-sensitive augmentation effects on
CIFAR-10, consistent with the positive OLS slope we observe; however, their
finding was that benefit decreases at large $N$, not that benefit is uniformly
negative.
This suggests the 5-epoch budget is uniquely constraining in a way that
inverts the typical outcome.
\citet{kang2024game} demonstrated augmentation benefits with MFG Augment, a
Mean-Field Game path-interpolation method that synthesises continuous image
distribution transitions, which is substantially more sophisticated than the
simple geometric transforms used here; their study also employed substantially
longer training runs.

\paragraph{The crossover estimate is not actionable.}
The OLS-extrapolated $N^\ast \approx 6.45$ million is presented for
completeness.
Given the bootstrap 95\% CI of $[40{,}325;\, {\sim}10^9]$ and the non-monotone
trend, this estimate should not be interpreted as a reliable prediction.
Any practical guidance derived from it would be speculative.

\paragraph{Practical implication.}
Practitioners should not apply basic geometric augmentation in short-run training
pipelines ($\leq 5$ epochs) without verifying that training time is sufficient
for the CNN to overcome convergence overhead.
The benefit of augmentation is epoch-budget-dependent, not merely
dataset-size-dependent.

%% ── 6. Limitations ────────────────────────────────────────────────────────────
\section{Limitations}
\label{sec:limitations}

\begin{enumerate}[leftmargin=*, label=\arabic*.]
\item \textbf{Two seeds per cell.}
  With $n{=}2$ replicates per (size, condition) cell, formal significance testing
  of per-size effects is not possible.
  The Spearman $\rho$ and OLS regression are purely descriptive; the
  $\rho{=}0.714$ ($p{=}0.071$) trend does not reach $\alpha{=}0.05$
  significance with $n{=}7$ size levels.

\item \textbf{Highly constrained epoch budget.}
  Five epochs is an extremely restricted training protocol.
  Results are specific to this under-trained regime and must not be generalised
  to longer training runs, where augmentation is well-established to provide
  benefit.

\item \textbf{Non-monotone mid-range trend.}
  The dip at $N{=}5{,}000$--$10{,}000$ (mean $\bar\Delta$ worsening from
  $-4.77$ to $-6.55$~pp between $N{=}2{,}500$ and $N{=}5{,}000$) reduces
  confidence in the log-linear model as a complete description of the
  size-dependence ($R^2{=}0.613$).

\item \textbf{Unreliable crossover estimate.}
  $N^\ast{=}6.45$ million is highly extrapolated from the tested range
  (max $N{=}60{,}000$); the bootstrap 95\% CI spanning four orders of magnitude
  provides negligible practical predictive value.

\item \textbf{Single architecture and augmentation policy.}
  Results apply to the specific dual-conv architecture with BatchNorm and the
  basic flip+affine policy; they may not hold for larger models, architectures
  without batch normalisation, or more sophisticated augmentation methods
  (Mixup, CutOut, AutoAugment, RandAugment).

\item \textbf{Fashion-MNIST only.}
  The specific epoch budget at which augmentation transitions from harmful to
  helpful may differ for datasets with higher intra-class complexity (e.g.,
  CIFAR-10).

\item \textbf{Architecture deviation from brief.}
  The implemented architecture uses 256-unit FC head and two convolutions per
  block (vs.\ 128-unit FC and one convolution per block in the original study
  brief); this matches the experiment plan's configuration file
  (\texttt{fc\_units=256}, \texttt{conv\_filters=[32,32,64,64]}) and is
  consistent throughout all 28 runs.
\end{enumerate}

%% ── 7. Conclusion ─────────────────────────────────────────────────────────────
\section{Conclusion}
\label{sec:conclusion}

We conducted a fully controlled study of basic geometric augmentation across
seven training set sizes on Fashion-MNIST, training a small CNN for exactly
5 epochs per run.
Augmentation was not beneficial under the 5-epoch budget at any tested size,
with mean $\bar\Delta$ ranging from $-8.06$~pp at $N{=}500$ to $-2.81$~pp at
$N{=}60{,}000$.
Perfect seed concordance (7/7 sizes, 2/2 seeds) confirms this is a systematic
effect, not a stochastic artefact.

Augmentation harm decreases as training set size grows ($R^2{=}0.613$), but the
estimated crossover to positive benefit at $N^\ast \approx 6.45$ million is
unreliably extrapolated and should not be treated as an actionable threshold.

The result is explained by \emph{convergence overhead}: augmentation induces
harder per-step examples, slowing convergence in a way that is not overcome
within 5 epochs regardless of training set size.
This finding does not imply that augmentation is generally harmful; it implies
that augmentation benefit is epoch-budget-dependent.

Future work should vary the epoch budget systematically to identify the minimum
number of epochs at which basic geometric augmentation begins to provide net
benefit at each training set size, and should extend the analysis to
architectures without batch normalisation and to datasets with higher intra-class
variability \citep{karada2020empirical,kang2024game,mukhamediev2024state}.

%% ── Reproducibility ───────────────────────────────────────────────────────────
\appendix
\section*{Appendix: Reproducibility Statement}
\label{app:repro}

All code, configuration, data indices, and result files are provided in the
accompanying workspace.
The study can be reproduced as follows.

\begin{enumerate}[leftmargin=*, label=\arabic*.]
\item \textbf{Data preparation.}
  Run \texttt{python src/make\_subsets.py} to generate stratified subset indices
  and a class-balance check; requires
  \texttt{torchvision.datasets.FashionMNIST} auto-download to \texttt{data/}.
\item \textbf{Training (seed~42).}
  Run \texttt{python src/gpu\_train\_seed42.py} on a CUDA-enabled GPU.
  Outputs per-epoch results (70 rows) and a device log to \texttt{results/}.
\item \textbf{Training (seed~123).}
  Run \texttt{python src/gpu\_train\_seed123.py} similarly.
\item \textbf{Analysis and figures.}
  Run \texttt{python scripts/t4\_3\_produce\_outputs.py} to reproduce all
  figures, the results table, and the summary report.
  Analysis artefacts (regression JSON, bootstrap CSV, diagnostic text) are in
  \texttt{analysis/}.
\item \textbf{Configuration.}
  All hyperparameters are frozen in \texttt{config/config.json}:
  learning rate $10^{-3}$, batch size 64, 5 epochs, Adam optimiser,
  flip probability 0.5, affine degrees $\pm10^\circ$,
  translate $(0.1, 0.1)$,
  normalisation mean/std $(0.2860,\, 0.3530)$.
\item \textbf{Environment.}
  NVIDIA~L4 GPU, CUDA~13.0, Python~3.13.15, PyTorch~2.11.0+cu130.
  Full environment specification is in \texttt{logs/gpu/\_environment.txt}.
\end{enumerate}

Random seeds 42 and 123 are applied to \texttt{random}, NumPy, and PyTorch
(including \texttt{torch.cuda.manual\_seed\_all}) at the start of each
training run.
Regression and bootstrap diagnostics were produced by automated analysis agents;
the regression uses OLS via \texttt{scipy.stats.linregress} with
predictor $\log_2(N)$.

%% ── Bibliography ──────────────────────────────────────────────────────────────
\bibliographystyle{plainnat}
\bibliography{references}

\end{document}
